\section{Bài báo giải được gì, và để lại gì}
\label{sec:ch05-con-lai-gi}

\index{encoder-decoder!bài giải được gì}%
Cái bài này giải được, phát biểu cho gọn: nó cho một mạng nơ-ron duy nhất, huấn
luyện đầu cuối, không giả định gì về cấu trúc bài toán, thắng một hệ dịch thống
kê trên một tác vụ dịch quy mô lớn. BLEU 34.81 cho ensemble năm LSTM đảo với
beam 12, so với 33.30 của hệ đối chứng. Bài nói rõ đây là lần đầu một hệ dịch
thuần nơ-ron thắng một hệ đối chứng cụm từ với khoảng cách đáng kể, và thắng dù
không xử lý nổi từ ngoài từ điển; nguyên văn là \enquote{by a sizeable margin,
despite its inability to handle out-of-vocabulary words}.

Cạnh con số ấy phải đặt thêm một con số nữa, không thì dễ đọc nhầm: 37.0, kết
quả WMT'14 tốt nhất lúc bấy giờ. Mạng thuần chưa vượt nó. Chỗ mạng chạm tới 36.5
là khi dùng để chấm điểm lại một nghìn giả thuyết của chính hệ đối chứng, tức
vẫn dựa lưng vào hệ thống cũ.

\subsection*{Cái bài tự nêu là còn thiếu}

\index{encoder-decoder!giới hạn bài tự nêu}%
Danh sách này ngắn, và ngắn theo cách đáng chú ý.

Từ ngoài \tn{từ điển}{vocabulary}. Từ điển đích 80 nghìn từ, và mọi từ ngoài đó
thành \code{UNK}; bài nói thẳng điểm BLEU của nó bị phạt mỗi lần bản dịch tham
chiếu chứa một từ như thế. Đây là giới hạn duy nhất bài nhắc lại nhiều lần.

Chi phí huấn luyện. Một GPU chạy khoảng 1700 từ mỗi giây, quá chậm để dùng, nên
họ chia mô hình ra tám GPU và đạt 6300 từ mỗi giây; tổng cộng khoảng mười ngày.
Bài trình bày đây là chuyện kỹ thuật đã giải xong chứ không phải một giới hạn.

Và hết. Bài \emph{không} nêu \tn{vector ngữ cảnh}{context vector} cố định là
một giới hạn, vì theo đo đạc của chính họ nó không phải giới hạn ở quy mô họ
chạy.

\subsection*{Cái chỉ nhìn từ chỗ khác mới thấy}

\index{vector ngữ cảnh!vì sao bài không thấy chỗ thắt}%
Đây là chỗ chương này phải cẩn thận nhất, vì hai câu đúng nằm cạnh nhau và
nghe như mâu thuẫn.

Câu thứ nhất: chỗ thắt là có thật.
Mục~\ref{sec:ch05-bop-vector} đo được nó, và đo được cả hình dạng của nó, tức
câu ngắn bão hòa trước còn câu dài cần vector rộng hơn hẳn. Ở \(d_h = 64\), câu
ngắn được 0.9797 còn câu dài được 0.0724.

Câu thứ hai: bài này không chạm vào nó. Mô hình của họ mang 8000 số thực cho một
câu. Bảng của tôi dừng ở 256, và ngay tại 256 thì câu dài đã lên 0.7105 rồi.

Hai câu ấy không đá nhau, và chỗ chúng khớp vào nhau là chỗ đáng nhớ:
\textbf{chỗ thắt không phải một tính chất của kiến trúc, nó là một tỷ lệ giữa bề
rộng vector và lượng thông tin phải đi qua nó.} Cùng một kiến trúc, thu hẹp
vector lại thì hỏng, nới ra thì không.
Bài chạy ở phía không hỏng và báo cáo đúng thứ họ thấy. Cho và cộng sự trong bài
2014b chạy ở phía kia và báo cáo đúng thứ họ thấy. Hai kết quả không đá nhau.

Nên cách nói \enquote{encoder-decoder hỏng ở câu dài} thiếu mất vế quan trọng
nhất, là vế \enquote{với vector rộng bao nhiêu}. Và cách sửa mà
chương~\ref{ch:dong-chinh} sẽ đọc không phải là nới vector ra, vì nới ra thì
lại hỏng ở câu dài hơn nữa. Nó là bỏ hẳn cái ràng buộc bề rộng cố định đi.

\subsection*{Ba thứ chương này để lại cho chương sau}

\index{vector ngữ cảnh!ba câu hỏi để ngỏ}%
Thứ nhất, và là thứ chính: encoder nén \tn{câu nguồn}{source sentence} thành
một vector rồi vứt bỏ
mọi trạng thái trung gian nó đã đi qua. Encoder \emph{đã} tính \(T\) trạng thái,
mỗi trạng thái một bước, và dùng đúng cái cuối cùng. Câu hỏi tự nhiên là vì sao
phải vứt \(T - 1\) cái kia, và đó là câu hỏi chương~\ref{ch:dong-chinh} mở ra.

Thứ hai, decoder không có cách nào hỏi lại. Nó viết từ thứ mười lăm dựa trên một
vector đã cố định từ trước khi nó viết từ nào, cộng với những gì chính nó vừa
viết. Không có bước nào trong \eqref{eq:ch05-phan-ra} nhìn lại câu nguồn.

Thứ ba, thứ mà chương~\ref{ch:lstm} để lại thì chương này không đụng tới. Cả
encoder lẫn decoder vẫn chạy từng bước một, và một câu \(T\) token vẫn là \(T\)
lượt tính nối đuôi nhau. Bài còn làm nó nặng thêm: bốn tầng, nên bốn lượt mỗi
bước. Mười ngày trên tám GPU là cái giá của chuyện đó ở quy mô sản xuất.

Chương~\ref{ch:dong-chinh} giải quyết hai thứ đầu và không đụng tới thứ ba. Nó
thêm một đường đi từ mọi bước của encoder tới mọi bước của decoder, đo cái chỗ
thắt \emph{trước} khi vá nó, rồi cho thấy đường mới ấy làm gì với gradient. Cái
giá là attention lúc đó vẫn là một thứ lắp thêm vào một mạng hồi quy, và mạng
hồi quy thì vẫn tuần tự. Phải tới chương~\ref{ch:transformer} mới có người
quyết định trả giá để bỏ hẳn phép hồi quy.
